%%
%% This is file `sample-sigconf-authordraft.tex',
%% generated with the docstrip utility.
%%
%% The original source files were:
%%
%% samples.dtx  (with options: `all,proceedings,bibtex,authordraft')
%% 
%% IMPORTANT NOTICE:
%% 
%% For the copyright see the source file.
%% 
%% Any modified versions of this file must be renamed
%% with new filenames distinct from sample-sigconf-authordraft.tex.
%% 
%% For distribution of the original source see the terms
%% for copying and modification in the file samples.dtx.
%% 
%% This generated file may be distributed as long as the
%% original source files, as listed above, are part of the
%% same distribution. (The sources need not necessarily be
%% in the same archive or directory.)
%%
%%
%% Commands for TeXCount
%TC:macro \cite [option:text,text]
%TC:macro \citep [option:text,text]
%TC:macro \citet [option:text,text]
%TC:envir table 0 1
%TC:envir table* 0 1
%TC:envir tabular [ignore] word
%TC:envir displaymath 0 word
%TC:envir math 0 word
%TC:envir comment 0 0
%%
%% The first command in your LaTeX source must be the \documentclass
%% command.
%%
%% For submission and review of your manuscript please change the
%% command to \documentclass[manuscript, screen, review]{acmart}.
%%
%% When submitting camera ready or to TAPS, please change the command
%% to \documentclass[sigconf]{acmart} or whichever template is required
%% for your publication.
%%
%%
% \documentclass[sigconf,authordraft]{acmart}
% \documentclass[manuscript, screen, review]{acmart}
\documentclass[sigconf,review,anonymous]{acmart}
%%
%% \BibTeX command to typeset BibTeX logo in the docs
\AtBeginDocument{%
  \providecommand\BibTeX{{%
    Bib\TeX}}}

%% Rights management information.  This information is sent to you
%% when you complete the rights form.  These commands have SAMPLE
%% values in them; it is your responsibility as an author to replace
%% the commands and values with those provided to you when you
%% complete the rights form.
\setcopyright{acmlicensed}
\copyrightyear{2018}
\acmYear{2018}
\acmDOI{XXXXXXX.XXXXXXX}
%% These commands are for a PROCEEDINGS abstract or paper.
\acmConference[Conference acronym 'XX]{Make sure to enter the correct
  conference title from your rights confirmation email}{June 03--05,
  2018}{Woodstock, NY}
%%
%%  Uncomment \acmBooktitle if the title of the proceedings is different
%%  from ``Proceedings of ...''!
%%
%%\acmBooktitle{Woodstock '18: ACM Symposium on Neural Gaze Detection,
%%  June 03--05, 2018, Woodstock, NY}
\acmISBN{978-1-4503-XXXX-X/2018/06}


%%
%% Submission ID.
%% Use this when submitting an article to a sponsored event. You'll
%% receive a unique submission ID from the organizers
%% of the event, and this ID should be used as the parameter to this command.
%%\acmSubmissionID{123-A56-BU3}

%%
%% For managing citations, it is recommended to use bibliography
%% files in BibTeX format.
%%
%% You can then either use BibTeX with the ACM-Reference-Format style,
%% or BibLaTeX with the acmnumeric or acmauthoryear sytles, that include
%% support for advanced citation of software artefact from the
%% biblatex-software package, also separately available on CTAN.
%%
%% Look at the sample-*-biblatex.tex files for templates showcasing
%% the biblatex styles.
%%

%%
%% The majority of ACM publications use numbered citations and
%% references.  The command \citestyle{authoryear} switches to the
%% "author year" style.
%%
%% If you are preparing content for an event
%% sponsored by ACM SIGGRAPH, you must use the "author year" style of
%% citations and references.
%% Uncommenting
%% the next command will enable that style.
%%\citestyle{acmauthoryear}


%%
%% end of the preamble, start of the body of the document source.
\begin{document}

%%
%% The "title" command has an optional parameter,
%% allowing the author to define a "short title" to be used in page headers.
\title{Preserving Human Agency in Agentic LLM Pitch Refinement Through Control-Preserving Cues}

%%
%% The "author" command and its associated commands are used to define
%% the authors and their affiliations.
%% Of note is the shared affiliation of the first two authors, and the
%% "authornote" and "authornotemark" commands
%% used to denote shared contribution to the research.
\author{Kudzai Sauka}
\email{k.sauka@hva.nl}
\orcid{ORCID: 0000-0002-3233-895X}
% \authornotemark[1]
\affiliation{%
  \institution{Amsterdam Univ. of Applied Sciences}
  \city{Amsterdam}
  \country{The Netherlands}
}

\author{Frederik Situmeang}
\email{f.b.i.situmeang@hva.nl}
\orcid{ORCID: 0000-0002-2156-2083}
% \authornotemark[1]
\affiliation{%
  \institution{Amsterdam Univ. of Applied Sciences}
  \city{Amsterdam}
  \country{The Netherlands}
}

\author{Monika Kackovic}
\email{m.kackovic@uva.nl}
\orcid{ORCID: 0000-0002-7423-3902}
% \authornotemark[1]
\affiliation{%
  \institution{Univ. of Amsterdam}
  \city{Amsterdam}
  \country{The Netherlands}
}


%%
%% By default, the full list of authors will be used in the page
%% headers. Often, this list is too long, and will overlap
%% other information printed in the page headers. This command allows
%% the author to define a more concise list
%% of authors' names for this purpose.
\renewcommand{\shortauthors}{Sauka et al.}

%%
%% The abstract is a short summary of the work to be presented in the
%% article.
\begin{abstract}
Large Language Model (LLM) assistants increasingly do more than respond to prompts. They propose next steps, execute subtasks, and shape task progress with limited user direction. This shift creates a shared-autonomy problem in Human-AI interaction, because the same system that improves efficiency may also reduce the user’s experience of control and authorship. This risk becomes sharper when the assistant’s actions are both explainable and anthropomorphic. Provenance-based explainability can make agentic actions appear intelligible and legitimate, while anthropomorphic communication cues can make those same actions easier to accept without resistance. The present study examines whether these cues increase users’ overreliance on AI advice, and whether overreliance reduces sense of agency. It further tests whether interrogative agendas preserve human agency by weakening the extent to which overreliance translates into agency loss. Interrogative agendas are implemented through questions, options, checkpoints, divergence signalling, and explicit decision-right reminders before the assistant advances the task. The study is conducted in an entrepreneurial pitch-refinement context, where efficiency gains matter but authorship and decision ownership must remain with the user. The paper contributes a human-centred account of how explainability and anthropomorphic communication cues may jointly deepen the agency problem in agentic LLM interaction, and introduces interrogative agendas as a design layer for preserving sense of agency during collaborative task progress.
\end{abstract}


%%
%% The code below is generated by the tool at http://dl.acm.org/ccs.cfm.
%% Please copy and paste the code instead of the example below.
%%
\begin{CCSXML}
<ccs2012>
 <concept>
  <concept_id>00000000.0000000.0000000</concept_id>
  <concept_desc>Do Not Use This Code, Generate the Correct Terms for Your Paper</concept_desc>
  <concept_significance>500</concept_significance>
 </concept>
 <concept>
  <concept_id>00000000.00000000.00000000</concept_id>
  <concept_desc>Do Not Use This Code, Generate the Correct Terms for Your Paper</concept_desc>
  <concept_significance>300</concept_significance>
 </concept>
 <concept>
  <concept_id>00000000.00000000.00000000</concept_id>
  <concept_desc>Do Not Use This Code, Generate the Correct Terms for Your Paper</concept_desc>
  <concept_significance>100</concept_significance>
 </concept>
 <concept>
  <concept_id>00000000.00000000.00000000</concept_id>
  <concept_desc>Do Not Use This Code, Generate the Correct Terms for Your Paper</concept_desc>
  <concept_significance>100</concept_significance>
 </concept>
</ccs2012>
\end{CCSXML}

\ccsdesc[500]{Do Not Use This Code~Generate the Correct Terms for Your Paper}
\ccsdesc[300]{Do Not Use This Code~Generate the Correct Terms for Your Paper}
\ccsdesc{Do Not Use This Code~Generate the Correct Terms for Your Paper}
\ccsdesc[100]{Do Not Use This Code~Generate the Correct Terms for Your Paper}

%%
%% Keywords. The author(s) should pick words that accurately describe
%% the work being presented. Separate the keywords with commas.
\keywords{agentic LLM, sense of agency, overreliance on AI advice, explainability, anthropomorphic communication cues, interrogative agendas}
%% A "teaser" image appears between the author and affiliation
%% information and the body of the document, and typically spans the
%% page.
% \begin{teaserfigure}
%   \includegraphics[width=\textwidth]{sampleteaser}
%   \caption{Seattle Mariners at Spring Training, 2010.}
%   \Description{Enjoying the baseball game from the third-base
%   seats. Ichiro Suzuki preparing to bat.}
%   \label{fig:teaser}
% \end{teaserfigure}

% \received{20 February 2007}
% \received[revised]{12 March 2009}
% \received[accepted]{5 June 2009}

%%
%% This command processes the author and affiliation and title
%% information and builds the first part of the formatted document.
\maketitle

\section{Introduction}

Large Language Model (LLM) assistants increasingly do more than answer prompts. They propose next steps, execute subtasks, and move tasks forward with limited user direction. This shift matters for Human-AI interaction because the central question is no longer only whether the system produces a useful output. It is also whether the user remains in control when the system begins to act with initiative. As LLM assistants become more agentic, interaction moves from simple turn-taking toward shared autonomy, where control, authorship, and decision ownership can become less stable during task progress.

This problem becomes sharper when agent initiative is supported by cues that make it easier to accept. Provenance-based explainability can make the assistant’s actions appear intelligible, justified, and legitimate, because the system does not only act but also provides a visible reason for acting. Anthropomorphic communication cues can add a second layer by making the same initiative easier to accept as socially appropriate. When these cues appear together, users may become more willing to follow agentic recommendations without pausing to evaluate whether those recommendations still reflect their own intent. The issue is therefore not only whether explainability or anthropomorphic communication improves interaction in general terms. The issue is whether these cues deepen overreliance on AI advice in ways that reduce sense of agency.

Entrepreneurial pitch refinement provides a useful context for examining this problem because it combines explicit task constraints with strong expectations of authorship. A founder may benefit from assistance that restructures a pitch, removes weak claims, and proposes stronger evidence, yet the same assistance may become uncomfortable if it begins to shape the task in ways that feel too persuasive or too far ahead of the founder’s own judgement. In this context, a better pitch does not necessarily imply a better interaction. Output quality may improve while sense of agency declines.

The present paper addresses this problem by examining whether interrogative agendas can reduce agency loss in agentic LLM interaction. Specifically, the study tests whether provenance-based explainability and anthropomorphic communication cues increase overreliance on AI advice, whether overreliance reduces sense of agency, and whether interrogative agendas weaken that negative relationship. In this study, interrogative agendas are defined as a mixed-initiative interaction design in which the system structures initiative through questions, options, checkpoints, and explicit decision-right reminders before advancing the task. Rather than preventing agent initiative, this design keeps initiative visible, contestable, and user-directed.

This paper makes three contributions. First, it develops a human-centred account of how provenance-based explainability and anthropomorphic communication cues may jointly deepen the agency problem in agentic interaction by increasing overreliance on the assistant’s actions. Second, it introduces interrogative agendas as a mixed-initiative interaction structure that can keep users in the loop without removing the practical value of agent initiative. Third, it provides an empirical framework for examining sense of agency as a central outcome in agentic LLM collaboration, rather than treating trust or satisfaction as sufficient indicators of successful interaction.

\section{Theoretical Background}

\subsection{Explainability and anthropomorphic communication cues}

The present model is anchored in the view that users process system behaviour through multiple cue types rather than through task utility alone. The Heuristic-Systematic Model is useful here because it distinguishes between cues that support deliberate judgement and cues that support more heuristic acceptance. In agentic interaction, provenance-based explainability can function as a systematic cue because it provides a visible basis for why the assistant proposes, revises, or advances a task action. This can increase the apparent legitimacy of agent initiative by reducing opacity around what the system is doing and why it is doing it. At the same time, anthropomorphic communication cues can function as heuristic social cues because they make the same initiative appear more socially appropriate, more intelligible as an interactional act, and less likely to be resisted. These two cue types are analytically distinct, but in agentic settings they may converge in their practical effect by increasing willingness to accept system action \cite{bucinca2021trust,vasconcelos2023explanations,epley2007seeinghuman,go2019humanizing}.

This distinction matters because explainability is often treated as a straightforward remedy for user concerns, and anthropomorphic communication is often treated as a surface feature that mainly affects trust or social presence. In the present paper, both are treated differently. The issue is not whether these cues make the system more likeable in general, but whether they make agent initiative easier to accept before the user has actively evaluated whether that initiative still reflects their own intent. Under those conditions, explainability and anthropomorphic communication may jointly deepen a control problem rather than simply improve interaction quality.

\subsection{Overreliance on AI advice}

The mediator in the present model is overreliance on AI advice. Prior work on human-AI decision-making shows that users often accept AI recommendations with insufficient verification, even when those recommendations are wrong, and that explanations do not automatically reduce that tendency \cite{bucinca2021trust}. More recent work shows that explanations can reduce overreliance only when the explanation lowers the practical cost of verification and makes engagement worthwhile \cite{vasconcelos2023explanations}. In this paper, overreliance is treated as the tendency to allow the assistant’s proposed or executed actions to proceed with reduced scrutiny, reduced verification, or reduced intervention. The construct is therefore narrower than trust and more directly tied to task progression under agentic assistance.

This is important because agentic systems alter interaction by changing who appears to be steering the task. Once a system proposes and advances subtasks on the user’s behalf, the relevant outcome is no longer only whether the user trusts the system or likes the output. The relevant outcome is also whether the user continues to experience authorship, intervention capacity, and decision ownership throughout the task. Overreliance is therefore treated here as the central mechanism through which explainability and anthropomorphic communication cues may reduce users’ sense of agency.

\subsection{Sense of agency under shared autonomy}

Sense of agency refers to the perceived control that individuals experience over their own actions and the consequences of those actions. In human-AI interaction, this construct has become increasingly important as systems move from passive tools to more autonomous collaborators. Recent work has shown that greater system autonomy can improve task performance while decreasing users’ experienced agency, especially in shared-autonomy settings where the system begins to shape action trajectories more actively \cite{legaspi2024sense,collier2025sense}. In the present study, sense of agency is treated as the degree to which users experience themselves as directing the task, owning key decisions, and retaining authorship over the final output.

This construct is narrower than broad control beliefs and more appropriate to the present task context, because the issue is not only whether users feel generally comfortable with the system but whether they still experience the refined pitch as their own decision product. The claim is therefore direct: higher overreliance on AI advice is expected to reduce users’ sense of agency.

\subsection{Mixed-initiative control cues and interrogative agendas}

If explainability and anthropomorphic communication cues can make agent initiative easier to accept, then a countervailing design layer is needed to preserve user control without removing the practical value of agentic assistance. Mixed-initiative interaction and human-AI design guidelines provide the literature background for this design. In those traditions, useful system initiative is not rejected, but it is structured so that users can see, question, redirect, or withhold acceptance before the system advances the task \cite{horvitz1999principles,amershi2019guidelines}. Interrogative agendas are grounded in the mixed-initiative interaction literature and are used in this study as the focal control-preserving design. Interrogative agendas are a form of control-preserving interaction in which the system structures agent initiative through questions, options, checkpoints, and explicit decision-right reminders before advancing the task. This matters because such cues do not eliminate system initiative. Instead, they slow silent acceptance, reinsert visible decision points, and preserve the user’s role as the directing agent in the interaction. This logic is also consistent with cognitive-forcing approaches that reduce overreliance by requiring more active engagement with AI outputs \cite{bucinca2021trust}. Accordingly, interrogative agendas are expected to weaken the negative relationship between overreliance on AI advice and sense of agency.

\section{Conceptual Model and Hypotheses}

Figure~\ref{fig:model} presents the conceptual model of agency loss and control preservation in agentic LLM interaction. Provenance-based explainability and anthropomorphic communication cues are expected to increase overreliance on AI advice. Overreliance is expected to reduce sense of agency. Interrogative agendas are expected to weaken the negative relationship between overreliance and sense of agency.

\begin{figure}[ht]
    \centering
    \fbox{\parbox{0.92\linewidth}{
    \centering
    \textbf{Provenance-Based Explainability} \hspace{1cm} \textbf{Anthropomorphic Communication Cues}\\[0.8em]
    $\downarrow$ \hspace{3.9cm} $\downarrow$\\[0.8em]
    \textbf{Overreliance on AI Advice}\\[0.8em]
    $\downarrow$ \, (negative effect on outcome)\\[0.8em]
    \textbf{Sense of Agency}\\[1em]
    \textit{Moderator: Interrogative Agendas}\\
    \textit{Covariates: Trust Propensity, Prior LLM Familiarity, Entrepreneurial Experience}
    }}
    \caption{Conceptual model of agency loss and control preservation in agentic LLM interaction.}
    \label{fig:model}
\end{figure}

Provenance-based explainability is expected to increase overreliance because visible rationale can make agentic action appear more intelligible and legitimate. Anthropomorphic communication cues are expected to increase overreliance because socially appropriate initiative is less likely to be resisted or interrupted. When these two cues are combined, their joint effect is expected to produce higher overreliance than either cue alone. Overreliance is then expected to reduce sense of agency by lowering intervention, resistance, and active authorship during task progression. Interrogative agendas are expected to preserve agency by weakening the extent to which overreliance translates into agency loss.

\begin{description}[leftmargin=0em,itemsep=0.4em]
    \item[\textbf{H1.}] Provenance-based explainability increases overreliance on AI advice.
    \item[\textbf{H2.}] Anthropomorphic communication cues increase overreliance on AI advice.
    \item[\textbf{H3.}] The combination of provenance-based explainability and anthropomorphic communication cues produces higher overreliance on AI advice than either cue alone.
    \item[\textbf{H4.}] Overreliance on AI advice is negatively associated with sense of agency.
    \item[\textbf{H5.}] Interrogative agendas weaken the negative relationship between overreliance on AI advice and sense of agency.
\end{description}

\section{Method}

\subsection{Study context and task}

The study is situated in entrepreneurial pitch refinement, a task context in which efficiency gains from system assistance are valuable, but authorship and decision ownership remain central. Participants are asked to refine an entrepreneurial pitch under explicit constraints, including concision, credibility, and investor relevance, while interacting with an agentic LLM assistant that can propose and advance revisions.

\subsection{Design}

The study manipulates three elements of the interaction. First, provenance-based explainability is varied to control whether the assistant provides visible rationale for proposed actions and revisions. Second, anthropomorphic communication cues are varied to control whether the assistant’s initiative is expressed in a more socially appropriate and human-like interactional form. Third, interrogative agendas are varied as the control-preserving interaction design, implemented as questions, options, checkpoints, divergence signalling, and explicit decision-right reminders before task advancement. The design is intended to test the antecedent effects of explainability and anthropomorphic communication on overreliance, and the moderating role of interrogative agendas in preserving users’ sense of agency.

\subsection{Measures}

\paragraph{Sense of Agency.}
Sense of Agency was measured as the degree to which participants experienced themselves as directing the task, owning key decisions, and retaining authorship over the final output during interaction with the assistant. This construct was adapted from the broader sense-of-agency literature, where sense of agency is defined as perceived control over one’s actions and their consequences, and from recent human-AI and shared-autonomy work showing that higher system autonomy can improve task performance while reducing users’ experienced agency \cite{tapal2017measure,legaspi2024sense,collier2025sense}. Because the present study focuses on entrepreneurial pitch refinement rather than motor control or generic autonomy judgements, the items were written to capture decision authorship, task direction, intervention capacity, and ownership of the final pitch.

Participants responded on a seven-point Likert scale (1 = strongly disagree, 7 = strongly agree). The scale items were: (a) I felt that I was directing the overall course of the task, (b) the important decisions in this task were still made by me, (c) I felt able to interrupt or redirect the assistant when needed, (d) the final pitch reflected my own judgement rather than only the assistant's suggestions, (e) I felt that I remained the author of the final output, and (f) during the task, I felt in control of how the work progressed.

\paragraph{Overreliance on AI Advice.}
Overreliance on AI Advice was measured as the tendency to allow the assistant's proposed or executed actions to proceed with reduced scrutiny, reduced verification, or reduced intervention. The measure was developed for the present study by drawing on adjacent literature on overreliance on AI, reliance on algorithmic advice, and verification behaviour in human-AI decision-making \cite{bucinca2021trust,vasconcelos2023explanations,schemmer2022meta}. In the present paper, overreliance is treated as a decision-process construct rather than as trust alone.

Participants responded on a seven-point Likert scale (1 = strongly disagree, 7 = strongly agree). The self-report items were: (a) I tended to accept the assistant's suggestions without examining them closely, (b) I often allowed the assistant to move the task forward without checking whether I agreed fully, (c) I accepted revisions from the assistant with little need to question them, (d) I rarely felt the need to stop and verify the assistant's proposed changes, (e) I let the assistant's suggestions shape the task even when I had not fully evaluated them, and (f) I found myself going along with the assistant's actions rather than actively directing each step.

\paragraph{Behavioural indicators of Overreliance on AI Advice.}
Because the overreliance literature shows that self-reported attitudes and actual reliance behaviour do not always align, overreliance was also captured behaviourally \cite{bucinca2021trust,vasconcelos2023explanations}. Four indicators were logged during the interaction: (1) the proportion of assistant-proposed revisions that were accepted, (2) the proportion of assistant-initiated actions accepted without user modification, (3) the number of user-initiated verification or clarification requests, and (4) the frequency of user-initiated interruptions, reversals, or redirections of assistant actions.

\paragraph{Secondary outcomes and covariates.}
Secondary outcomes include trust, perceived AI agency, autonomy comfort, and delegation. Trust propensity, prior LLM familiarity, and entrepreneurial experience are included as covariates.

\subsection{Analysis strategy}

The main analysis tests the effects of provenance-based explainability and anthropomorphic communication cues on overreliance on AI advice, the effect of overreliance on sense of agency, and the moderating role of interrogative agendas on the overreliance--agency relationship. Interaction terms are used to test the combined effect of explainability and anthropomorphic communication cues on overreliance, and the buffering effect of interrogative agendas on agency loss. Secondary outcomes are analysed separately and interpreted as complementary rather than central to the paper’s contribution.

\section{Results}

[Results to be inserted following data analysis. The reporting structure should move in the same order as the conceptual model: effects of explainability and anthropomorphic communication cues on overreliance, interaction effect on overreliance, effect of overreliance on sense of agency, and moderation by mixed-initiative control cues. Trust, perceived AI agency, autonomy comfort, and delegation should be reported as secondary outcomes.]

\section{Discussion}

The present paper treats explainability and anthropomorphic communication not as uniformly beneficial design features, but as cues that may jointly deepen the agency problem in agentic interaction by increasing overreliance on AI advice. If the hypothesised pattern is supported, the findings will indicate that users can remain comfortable, trusting, or even satisfied while still experiencing a reduction in sense of agency. That distinction matters because it shows that successful Human-AI interaction in agentic settings cannot be inferred from acceptance-oriented outcomes alone.

The paper also positions interrogative agendas as a practical intervention for this problem. Interrogative agendas do not remove the value of agent initiative. Instead, they structure initiative so that decision rights remain visible and active while the task progresses. This makes them suitable for settings such as pitch refinement, where support is useful but authorship cannot be ceded without changing the meaning of the task itself. More broadly, the paper suggests that the design of agentic assistants should be evaluated not only in terms of usefulness, trust, or fluency, but also in terms of whether user control remains visible and contestable throughout interaction.

\section{Limitations and Future Research}

The study is bounded by a specific task context, namely entrepreneurial pitch refinement, and by a controlled interaction setting. This improves conceptual discipline but limits immediate generalisation to longer-horizon tasks where authorship, iteration, and delegation unfold over extended periods. Future research should test whether the same mechanism operates in other agentic contexts, including collaborative planning, venture decision support, and more sustained entrepreneurial work. Future research should also examine whether different forms of explainability and different anthropomorphic communication cues produce distinct overreliance pathways rather than the more aggregated cue effects tested here.

\section{Conclusion}

As LLM assistants become more agentic, the central Human-AI interaction problem is no longer only whether the system can act effectively. It is also whether users remain in control when the system’s actions are both explainable and easy to accept. The present paper develops a model in which provenance-based explainability and anthropomorphic communication cues increase overreliance on AI advice, overreliance reduces sense of agency, and mixed-initiative control cues preserve human agency by weakening that translation. In this way, the paper positions interrogative agendas as a practical design form for keeping users in the loop while preserving the benefits of agentic assistance.

\begin{acks}
[Removed for anonymous review.]
\end{acks}

%%
%% The next two lines define the bibliography style to be used, and
%% the bibliography file.
\bibliographystyle{ACM-Reference-Format}
\bibliography{hai_acm_references_v2}


%%
%% If your work has an appendix, this is the place to put it.
\appendix


\end{document}
\endinput
%%
%% End of file `sample-sigconf-authordraft.tex'.
